\section{Innledning}

Trådløs kommunikasjon bygger på overføring og mottak av elektromagnetiske signaler ved bestemte frekvenser. For å forstå hvordan slike systemer fungerer, er det nyttig å undersøke hvordan en elektronisk krets kan generere høyfrekvente signaler, og hvordan komponentvalg, kretsutforming og fysisk montering påvirker signalets frekvens og styrke.
\\[1em]
I dette prosjektet er det arbeidet med en RF-krets som skal generere et signal i nærheten av FM-båndet. Kretsen består blant annet av en NE555-timer, en transistor, en LC-resonanskrets og en antenne. Gjennom teoretiske beregninger, praktisk bygging og målinger undersøkes sammenhengen mellom kretsens oppbygning og signalet som genereres.
\\[1em]
Hensikten med rapporten er å analysere hvordan kretsen er bygget opp, beregne forventede verdier for sentrale deler av kretsen og gjennomgå grundig feilsøking med målinger. Rapporten skal også vurdere hvilke faktorer som påvirker om kretsen oppnår ønsket funksjon.